\clearpage
\SourceLeaf{019}\pdfbookmark[0]{Book One.}{liber-primus}
\begin{center}
{\Large BOOK ONE.}\par\bigskip
{\large GENERAL RELATIONS AMONG THE QUANTITIES THAT DETERMINE THE MOTIONS OF CELESTIAL BODIES AROUND THE SUN.}\par\bigskip
\pdfbookmark[1]{Section One.}{sectio-prima}
{\large SECTION ONE.}\par\medskip
\emph{Relations concerning a single position in the orbit.}\par\bigskip
\pdfbookmark[2]{1.}{article-001}1.
\end{center}
In this work we shall consider the motions of celestial bodies only insofar as they are governed by the Sun's attractive force. All secondary planets are therefore excluded from our undertaking, as are the perturbations that the primary planets exert on one another and all rotational motion. We regard the moving bodies themselves as mathematical points and assume that all motions take place according to the following laws, which are therefore to be taken as the basis of all the investigations in this work.

I. The motion of any celestial body always takes place in one and the same plane, which also contains the centre of the Sun.

II. The trajectory described by the body is a conic section with one focus at the centre of the Sun.

III. Motion along this trajectory takes place in such a way that the areas swept out around the Sun in different
\SourceLeaf{020}intervals of time are proportional to those intervals. Thus, when the times and areas are expressed numerically, any area divided by the time in which it is swept out gives a constant quotient.

IV. For different bodies moving around the Sun, the squares of these quotients are in the compound ratio of the parameters corresponding to their orbits and of the respective sums of the Sun's mass and the masses of the moving bodies.

Let $2p$ denote the parameter of the orbit in which the body moves, $\mu$ the quantity of matter in that body (taking the Sun's mass as $=1$), and $\frac{1}{2}g$ the area it sweeps out around the Sun in time $t$. Then $\frac{g}{t\BareRoot{p}\cdot\BareRoot{(1+\mu)}}$ has the same value for all celestial bodies. Since it makes no difference which body we use to determine the value of this number, we shall derive it from the motion of the Earth, whose mean distance from the Sun we shall adopt as the unit of distance. Our unit of time will always be the mean solar day. Further, denoting by $\pi$ the ratio of a circle's circumference to its diameter, the area of the entire ellipse described by the Earth is evidently $\pi\BareRoot{p}$. This must therefore be set $=\frac{1}{2}g$ if the sidereal year is taken for $t$, in which case our constant becomes~$=\frac{2\pi}{t\BareRoot{(1+\mu)}}$. To find the numerical value of this constant, which we shall denote in what follows by $k$, we shall take from the latest determination the sidereal year, or $t=365{,}2563835$, and the mass of the Earth, or $\mu=\frac{1}{354710}=0{,}000\,0028\,192$. This gives

\begin{center}
\begin{tabular}{@{}ll@{}}
\makebox[15em][l]{$\log 2\pi$\hfill .\,.\,.\,.\,.\,.\,.\,.\,.\,.\,.\,.\,.\,.} & $0{,}798\,1798\,684$\\[.3em]
\makebox[15em][l]{Compl. $\log t$\hfill .\,.\,.\,.\,.\,.\,.\,.\,.\,.} & $7{,}437\,4021\,852$\\[.3em]
\makebox[15em][l]{Compl. $\log\BareRoot{(1+\mu)}$\hfill .\,.\,.\,.\,.} & $9{,}999\,9993\,878$\\[.3em]
\hline
\makebox[15em][l]{$\log k$\hfill .\,.\,.\,.\,.\,.\,.\,.\,.\,.\,.\,.\,.\,.\,.} & $8{,}235\,5814\,414$\\[.3em]
\multicolumn{1}{c}{$k=$} & $0{,}017\,2020\,9895$
\end{tabular}
\end{center}
